% Generated semantic source for AI Almanach v3.
% Edit v3 directly after initial generation; do not overwrite
% editorial changes by rerunning the converter casually.

\chapter{Practical LLM Engineering}
\chapterlead{Which mechanisms and trade-offs appear in a small LLM experiment?}{Trace}{Adapt}{Benchmark}{Deploy}{practical-llm-engineering}

\label{ch:practical-llm-engineering}

\begin{learningobjectives}
After this chapter, you should be able to:
\begin{itemize}
    \item trace text through tokenisation, embeddings, attention, MLP blocks,
          and logits;
    \item calculate the shapes and parameter counts of important Transformer
          matrices;
    \item derive the LoRA update and explain which parameters receive gradients;
    \item design a small supervised post-training dataset without leaking
          evaluation cases;
    \item compare a base model, an instruct model, and several adapter
          configurations;
    \item measure time to first token, decode throughput, memory, and output
          quality;
    \item decide when to keep an adapter separate and when to merge it into the
          base model;
    \item produce a reproducible evaluation dossier rather than a persuasive
          demo.
\end{itemize}
\end{learningobjectives}

\begin{plainlanguage}
An LLM is a repeated stack of matrix operations that predicts a distribution
for the next token.
Post-training does not pour facts into a tiny mechanical brain.
It changes selected numerical directions so that some continuations become
more likely and others less likely under a defined prompt format.
LoRA makes that change economical by learning two narrow matrices instead of
relearning each large weight matrix.
\end{plainlanguage}

\section{The Lab Contract}

This chapter uses \texttt{HuggingFaceTB/SmolLM2-135M-Instruct} as a concrete,
small decoder-only model.
The model card describes a 135-million-parameter member of the SmolLM2 family,
trained on two trillion tokens and then instruction-tuned with supervised
fine-tuning and preference
optimisation.\textsuperscript{\cite{benallalSmolLM2WhenSmol2025,
huggingfaceSmolLM2135MInstruct2026}}
The checkpoint is small enough for teaching and local experiments, but small
does not mean magically capable.

\begin{cautionbox}[Verify before running]
Model repositories, licenses, files, chat templates, and library APIs can
change.
Before an experiment, archive the model-card revision, license, exact commit or
file hashes, tokenizer configuration, and software versions.
The architectural values in \cref{tab:smollm2-config} were verified against the
published checkpoint configuration on 26 August 2026; they are not a promise
about a future revision.
\end{cautionbox}

\begin{table}[htbp]
\centering
\caption{Selected published configuration values for the teaching checkpoint.}
\label{tab:smollm2-config}
\begin{tblr}{
    width=\textwidth,
    colspec={X[l] Q[r,0.18\textwidth] X[l]},
    row{1}={font=\bfseries, bg=\maincolor!15!white},
    hlines,
    vlines}
Quantity & Value & Engineering consequence \\
Vocabulary size & 49,152 & Embedding and output spaces contain 49,152 token
IDs \\
Hidden size $d_{\mathrm{model}}$ & 576 & Each token state has 576 scalar
components \\
Transformer layers & 30 & Attention and MLP transformations repeat 30 times \\
Query heads & 9 & Each query head has dimension $576/9=64$ \\
Key/value heads & 3 & Three query heads share each KV head; the KV cache is
smaller \\
MLP intermediate size & 1,536 & Each token expands before being projected back
to 576 \\
Maximum configured positions & 8,192 & Longer inputs require a supported
strategy and fresh evaluation \\
Tied token embeddings & yes & Input embedding and output projection share
parameters \\
\end{tblr}
\end{table}

\begin{engineernote}
The lab is deliberately a model-comparison experiment.
Its deliverable is not ``the model said the desired sentence once.''
The deliverable is evidence showing what changed, what did not, what it cost,
and where it failed.
\end{engineernote}

\section{From a Prompt to the Next Token}

\factvisual{
    \textbf{One prediction is a pipeline.}
    The data shape changes at every boundary, but the model ultimately returns
    one score per vocabulary item.
    \begin{enumerate}
        \item format messages;
        \item tokenise text;
        \item look up vectors;
        \item mix positions with attention;
        \item transform each position with an MLP;
        \item map the final state to logits and probabilities.
    \end{enumerate}
}{
    \begin{tikzpicture}[node distance=3.2mm]
        \node[almanach card,text width=29mm] (text) {messages\\\scriptsize text};
        \node[almanach card,text width=29mm,below=of text] (ids)
            {tokeniser\\\scriptsize integer IDs};
        \node[almanach card,text width=29mm,below=of ids] (embed)
            {embedding\\\scriptsize $T\times d$ vectors};
        \node[almanach accent card,text width=29mm,below=of embed] (blocks)
            {Transformer blocks\\\scriptsize attention + MLP};
        \node[almanach card,text width=29mm,below=of blocks] (logits)
            {output projection\\\scriptsize $V$ logits};
        \node[almanach accent card,text width=29mm,below=of logits] (next)
            {decoding rule\\\scriptsize next token};
        \draw[almanach arrow] (text) -- (ids);
        \draw[almanach arrow] (ids) -- (embed);
        \draw[almanach arrow] (embed) -- (blocks);
        \draw[almanach arrow] (blocks) -- (logits);
        \draw[almanach arrow] (logits) -- (next);
    \end{tikzpicture}
    \visualcaption{The prompt-to-token path. Shapes and semantics change at
      every boundary; the next-token loop then appends the chosen token and
      repeats.}{fig:prompt-to-token-pipeline}
}

\subsection{Step 1: apply the chat template}

An instruct checkpoint expects role markers and control tokens in a particular
format.
The tokenizer's \keyterm{chat template} converts a list of system, user, and
assistant messages into that format.
Hand-written imitation templates can shift every token boundary and make an
otherwise correct training run teach the wrong interface.

\subsection{Step 2: tokenise}

The tokenizer maps text to integer IDs,

\begin{equation}
    (s_1,\ldots,s_m) \longmapsto (t_1,\ldots,t_T),
    \qquad t_i\in\{0,\ldots,V-1\},
    \label{eq:text-to-token-ids}
\end{equation}

where $V$ is vocabulary size and $T$ is sequence length.
Tokens are not necessarily words.
Whitespace, punctuation, code, and different languages can produce very
different token counts.

\begin{plainlanguage}[Tokenisation analogy---and its limit]
Tokenisation resembles replacing reusable text fragments with catalogue
numbers.
Unlike Morse code or a lossless compressor designed only for short bit strings,
the token inventory is also part of the model's learned computational interface.
Changing it changes the embedding and output spaces.
\end{plainlanguage}

\subsection{Step 3: look up embeddings}

Let the embedding matrix be $E\in\mathbb{R}^{V\times d}$.
Each token ID selects one row, producing

\begin{equation}
    X^{(0)} = E[t_1,\ldots,t_T]
    \in\mathbb{R}^{T\times d}.
    \label{eq:embedding-lookup}
\end{equation}
\where{$t_1,\dots,t_T$ are the token IDs, $E$ the embedding table, $T$ the sequence
length, and $d$ the model dimension; the operation is a table lookup, not a
matrix multiplication.}

For this checkpoint, $V=49{,}152$ and $d=576$, so the embedding table contains
$49{,}152\times576=28{,}311{,}552$ scalar parameters.
Because the output weights are tied, the same parameter table is reused when
hidden states are mapped back to token logits.

\subsection{Step 4: attention mixes information across positions}

For one conventional attention head with row-vector activations,

\begin{align}
    Q &= XW_Q, & K &= XW_K, & V &= XW_V,
    \label{eq:qkv-projections}\\
    A &= \operatorname{softmax}\!\left(
        \frac{QK^{\mathsf T}}{\sqrt{d_h}} + M\right),
    & H &= AV.
    \label{eq:scaled-dot-product-attention}
\end{align}
\where{$X$ is the layer input, $W_Q,W_K,W_V$ the learned projections, $d_h$ the
per-head dimension, and $M$ the causal mask that sets future positions to
$-\infty$ before the softmax.}

$M$ is a causal mask that prevents a position from reading future tokens.
The $T\times T$ matrix $A$ contains data-dependent mixing weights.
Multi-head attention performs related computations in several learned
projection spaces and combines their
outputs.\textsuperscript{\cite{vaswaniAttentionAllYou2023}}

The teaching checkpoint uses \keyterm{grouped-query attention}: nine query
heads share three key/value heads.
With head dimension 64, query projections produce $9\times64=576$ components,
while key and value projections each produce $3\times64=192$ components.
This reduces the key/value cache used during autoregressive generation.

\begin{engineernote}[Generation and the KV cache]
Autoregressive decoding still produces one new token per decoder step.
With a KV cache, earlier key and value tensors are reused; the model does not
need to recompute the complete prefix at every step.
Prompt processing and token-by-token decoding therefore have different cost
profiles and should be measured separately.
\end{engineernote}

\subsection{Step 5: the MLP transforms each position}

Attention exchanges information between positions.
The feed-forward block then applies learned transformations independently to
each position.
A Llama-style gated MLP can be written schematically as

\begin{equation}
    \operatorname{MLP}(x) =
    W_{\mathrm{down}}
    \left(\operatorname{SiLU}(W_{\mathrm{gate}}x)
    \odot W_{\mathrm{up}}x\right),
    \label{eq:gated-mlp}
\end{equation}

where $\odot$ is elementwise multiplication.
Residual connections add transformed information to the existing stream, and
normalisation controls scale.
The exact order and normalisation type are architectural choices; not every
Transformer uses the same pre-norm or post-norm design.

\subsection{Step 6: map hidden state to logits}

After the final block and normalisation, the last relevant hidden state is
projected to one logit per vocabulary item.

\begin{equation}
    z = hE^{\mathsf T}\in\mathbb{R}^{V},
    \qquad
    p(t_{T+1}=j\mid t_{\leq T})=
    \frac{\exp(z_j)}{\sum_{k=1}^{V}\exp(z_k)}.
    \label{eq:next-token-softmax}
\end{equation}

Greedy decoding chooses the largest probability.
Sampling transforms the distribution through temperature, top-$k$, top-$p$,
or related rules and then draws a token.
Decoding configuration is part of an evaluation protocol, not an aesthetic
afterthought.

\section{Where Approximately 135 Million Parameters Live}

Parameter counting is a dimensional-analysis exercise and an excellent way to
catch misunderstandings.
Ignoring small normalisation vectors, one layer of the published configuration
has approximately:

\begin{align}
    N_{\mathrm{attention}}
    &= 576\cdot576
      + 192\cdot576
      + 192\cdot576
      + 576\cdot576
      = 884{,}736,
      \label{eq:smollm-attention-parameters}\\
    N_{\mathrm{MLP}}
    &= 3\cdot(1{,}536\cdot576)
      = 2{,}654{,}208.
      \label{eq:smollm-mlp-parameters}
\end{align}
\where{the four terms are the query, key, value, and output projections; the key
and value terms are smaller because this model shares three query heads per
key/value head (grouped-query attention).}

Thirty layers therefore contribute about $106.17$ million parameters, and the
tied embedding table adds about $28.31$ million.
Normalisation vectors bring the total to roughly $134.5$ million, which is
appropriately described as a 135M model.

\begin{checkpoint}
Check the units in \cref{eq:smollm-attention-parameters}.
Why are the key and value projections smaller than the query and output
projections?
What would happen to embedding parameters if vocabulary size doubled while
hidden size stayed fixed?
\end{checkpoint}

\section{LoRA: Learning a Small Weight Update}

Full fine-tuning learns an update for every selected base parameter.
Low-Rank Adaptation freezes a base matrix and parameterises its update as the
product of two narrower matrices.\textsuperscript{\cite{huLoRALowRank2021}}

Deep-learning frameworks commonly store a linear-layer weight as
$W\in\mathbb{R}^{d_{\mathrm{out}}\times d_{\mathrm{in}}}$ and compute
$Y=XW^{\mathsf T}$.
LoRA defines

\begin{align}
    \Delta W &= \frac{\alpha}{r}BA,
    \label{eq:lora-update}\\
    A &\in\mathbb{R}^{r\times d_{\mathrm{in}}},
    &B &\in\mathbb{R}^{d_{\mathrm{out}}\times r},
    \label{eq:lora-shapes}\\
    Y' &= X(W+\Delta W)^{\mathsf T}.
    \label{eq:lora-forward}
\end{align}
\where{$r$ is the adapter rank, $\alpha$ a scaling constant, and $A,B$ the two
low-rank factors; only $A$ and $B$ receive gradients, while the base weights $W$
stay frozen.}

Because $\operatorname{rank}(BA)\leq r$, LoRA can only express updates in a
low-rank family, which is the deliberate constraint.

\begin{plainlanguage}[The two-narrow-matrix intuition]
Imagine that the original weight matrix is a large routing table.
Full fine-tuning may alter every route independently.
LoRA instead learns $r$ reusable input patterns in $A$ and $r$ corresponding
output directions in $B$.
Their combinations create a structured correction to the original table.
\end{plainlanguage}

The number of trainable LoRA scalars for one matrix is

\begin{equation}
    N_{\mathrm{LoRA}} = r(d_{\mathrm{in}}+d_{\mathrm{out}}),
    \label{eq:lora-parameter-count}
\end{equation}
\where{$d_{\mathrm{in}}$ and $d_{\mathrm{out}}$ are the dimensions of the weight
matrix being adapted. Compare $d_{\mathrm{in}}d_{\mathrm{out}}$ for the full
matrix: the saving is the whole point of the method.}

compared with $d_{\mathrm{in}}d_{\mathrm{out}}$ base scalars.

\begin{workedexample}[Rank-eight query and value adapters]
For the 576-by-576 query matrix and rank $r=8$,

\[
    N_Q = 8(576+576)=9{,}216.
\]

For the 192-by-576 value matrix,

\[
    N_V = 8(576+192)=6{,}144.
\]

Targeting query and value projections in all 30 layers therefore trains
$30(9{,}216+6{,}144)=460{,}800$ LoRA parameters, about \qty{0.34}{\percent} of
a 134.5-million-parameter base.
This estimate must be checked against the framework's printed trainable
parameter count because target-module names and architecture details can vary.
\end{workedexample}

\subsection{When gradients are computed}

During a training step:

\begin{enumerate}
    \item the batch is tokenised and padded or packed;
    \item a forward pass computes logits and causal language-modelling loss;
    \item automatic differentiation propagates gradients through the network;
    \item gradients with respect to frozen $W$ are not stored for optimisation;
    \item gradients for $A$ and $B$ are accumulated;
    \item the optimiser updates $A$ and $B$, then gradients are cleared.
\end{enumerate}

If $B$ begins at zero and $A$ begins randomly, then $\Delta W=0$ at
initialisation, so the adapter initially reproduces the base layer.
At the first step, $B$ can receive a gradient while the gradient for $A$ may be
zero; after $B$ moves away from zero, both factors can learn.

\begin{cautionbox}[LoRA does not make all training memory disappear]
Freezing base weights removes their optimiser states and parameter gradients,
but backpropagation still needs activations or recomputation to obtain adapter
gradients.
Sequence length, batch size, precision, checkpointing, and attention
implementation can still dominate memory.
\end{cautionbox}

\section{Supervised Post-Training Data}

For a token sequence $t_1,\ldots,t_T$, causal language modelling minimises the
negative log-likelihood of the next token,

\begin{equation}
    \mathcal{L}_{\mathrm{CLM}}
    = -\sum_{i=1}^{T-1}
      m_{i+1}\log p_{\theta}(t_{i+1}\mid t_{\leq i}),
    \label{eq:causal-lm-loss}
\end{equation}

where mask $m_{i+1}$ selects which tokens contribute to the loss.
An instruction-tuning pipeline may train on all formatted tokens or only on
assistant-response tokens.
Whichever rule is chosen must be documented and tested on one visible example.

\subsection{A useful example schema}

Store structured messages rather than one preformatted string when possible:

\begin{lstlisting}[language=json]
{"messages": [
  {"role": "system", "content": "Answer as a concise lab assistant."},
  {"role": "user", "content": "Convert 2.5 kW to watts."},
  {"role": "assistant", "content": "2.5 kW = 2500 W."}
]}
\end{lstlisting}

The experiment should use a licensed, narrow dataset whose desired behaviour
can be scored.
A handful of sentences is enough for a software smoke test, not enough to
support a broad claim of new capability.

\subsection{Split by the source of similarity}

Random row splits are weak when prompts are generated from templates.
Hold out entire prompt families, source documents, entities, or time periods so
that surface variants cannot straddle train and test.
Deduplicate exact and near-duplicate prompts and responses before splitting.
Keep a private challenge set whose prompts are never used in training,
debugging, early stopping, or prompt design.

\begin{table}[htbp]
\centering
\caption{Minimum dataset record for the tiny-model lab.}
\label{tab:tiny-llm-data-record}
\begin{tblr}{
    width=\textwidth,
    colspec={Q[l,0.24\textwidth] X[l]},
    row{1}={font=\bfseries, bg=\maincolor!15!white},
    hlines,
    vlines}
Field & Purpose \\
Stable example ID and hash & Detect duplicates and prove which example was
used \\
Source and license & Establish provenance and permitted use \\
Prompt-family or group ID & Enable leakage-resistant splitting \\
Messages before templating & Preserve semantic structure independently of
tokenizer version \\
Expected properties or rubric & Make evaluation possible without copying a
preferred phrase \\
Risk tags & Identify safety, privacy, formatting, and subgroup cases \\
\end{tblr}
\end{table}

\section{Training Conditions and Baselines}

One run cannot distinguish learning from randomness, leakage, prompt effects,
or a broken comparison.
Use a baseline matrix in which only declared factors change.

\begin{table}[htbp]
\centering
\caption{Baseline ladder for the LoRA experiment.}
\label{tab:lora-baselines}
\begin{tblr}{
    width=\textwidth,
    colspec={Q[c,0.08\textwidth] Q[l,0.25\textwidth] X[l]},
    row{1}={font=\bfseries, bg=\maincolor!15!white},
    hlines,
    vlines}
ID & Condition & Question answered \\
B0 & Released base checkpoint & What does pre-training alone provide? \\
B1 & Released instruct checkpoint & What did general instruction post-training
already provide? \\
B2 & Instruct checkpoint plus rank-4, rank-8, and rank-16 LoRA & Does rank
affect the quality--cost trade-off? \\
B3 & Attention-only versus attention-plus-MLP targets & Which modules need
adaptation? \\
B4 & No-op or randomly initialised adapter & Does the pipeline create an
apparent gain without learning? \\
B5 & Full fine-tuning, only if resources permit & What quality is traded for
parameter efficiency? \\
\end{tblr}
\end{table}

Keep dataset, splits, tokenizer, token budget, evaluation prompts, decoding,
and stopping rule fixed across candidate comparisons.
Use more than one seed for claims about training quality.
Record learning curves and the best checkpoint rule rather than reporting only
the winner.

\begin{engineernote}[A robust first-run checklist]
Before a long run, verify that the formatted input is readable, labels are
shifted correctly, masked positions are intentional, only adapter parameters
are trainable, loss is finite, one optimiser step changes adapter weights, a
saved adapter reloads, and a tiny batch can be overfit.
These tests catch more bugs per joule than heroic hyperparameter tuning.
\end{engineernote}

\section{Evaluation: Quality Is a Vector}

Language-model evaluation is sensitive to prompt templates, few-shot examples,
decoding, library versions, and task implementations.
Reproducible harnesses help, but do not make these choices
disappear.\textsuperscript{\cite{bidermanLessonsTrenchesReproducible2026}}
Holistic evaluation treats accuracy, calibration, robustness, fairness, bias,
toxicity, and efficiency as distinct dimensions rather than forcing them into
one leaderboard
number.\textsuperscript{\cite{liangHolisticEvaluationLanguage2023}}

\begin{table}[htbp]
\centering
\caption{Evaluation layers for the tiny-model experiment.}
\label{tab:tiny-llm-evaluation-layers}
\begin{tblr}{
    width=\textwidth,
    colspec={Q[l,0.20\textwidth] X[l] X[l]},
    row{1}={font=\bfseries, bg=\maincolor!15!white},
    hlines,
    vlines}
Layer & Example measures & Important control \\
Task correctness & Exact match, numeric tolerance, F1, schema validity &
Score held-out prompt families, not training paraphrases \\
Instruction following & Constraint satisfaction and rubric items &
Blind the evaluator to model identity and randomise order \\
Robustness & Paraphrase, whitespace, spelling, irrelevant context &
Preserve intended semantics and record perturbations \\
Safety and privacy & Refusal cases, sensitive-data regurgitation, prompt
injection &
Use authorised synthetic or public test data \\
General regression & A small unrelated capability suite &
Check that narrow gains did not destroy useful behaviour \\
Systems & Load time, TTFT, tokens/s, p50/p95 latency, memory &
Pin hardware, dtype, runtime, lengths, batch, and warm-up \\
\end{tblr}
\end{table}

\subsection{Avoid evaluator leakage}

Do not use the same model to generate training answers, design evaluation
rubrics, and judge outputs without independent checks.
For subjective tasks, use a blinded rubric with observable criteria, multiple
raters where possible, inter-rater agreement, and adjudication rules.
Show representative failures, not only averages.

\subsection{Report uncertainty}

For case-level metrics, use a paired bootstrap over independent prompt groups.
For human ratings, report the number of raters and cases, agreement, and an
interval suited to the sampling design.
For training comparisons, include variation across planned seeds.

\section{Benchmarking Latency, Throughput, and Memory}

\keyterm{Time to first token} measures the delay until the first generated
token and is strongly influenced by prompt processing.
\keyterm{Inter-token latency} or its reciprocal, decode tokens per second,
describes the autoregressive phase.
\keyterm{Throughput} may count output tokens across a batch or concurrent
requests and must therefore state its denominator.

\begin{practicallab}[Performance measurement matrix]
Compare base, separate-adapter, and merged-model inference under these fixed
conditions:
\begin{itemize}
    \item batch sizes 1 and 4;
    \item prompt lengths 32, 256, and 1,024 tokens;
    \item generated lengths 32 and 128 tokens;
    \item greedy decoding for deterministic timing;
    \item a pinned device, dtype, runtime, checkpoint revision, and power mode;
    \item at least five warm-up runs followed by at least 30 measured repeats.
\end{itemize}

Record checkpoint disk size, load time, host RAM, accelerator memory, time to
first token, output tokens per second, and end-to-end p50 and p95 latency.
Synchronise asynchronous accelerators before stopping a timer.
Do not compare numbers copied from different machines as if they were a
controlled experiment.
\end{practicallab}

\begin{cautionbox}[Why this chapter prints no universal tokens-per-second value]
Throughput depends on processor, memory bandwidth, kernels, batch size,
sequence lengths, dtype, quantisation, cache, thermal state, runtime version,
and timing method.
A number without those conditions is not portable evidence.
The lab teaches how to obtain a defensible number on the learner's system.
\end{cautionbox}

\section{Separate Adapter or Merged Model?}

PEFT-style tooling can keep LoRA factors separate or merge the update into a
standalone set of base weights.\textsuperscript{\cite{huggingfacePEFTLoRA2026}}
For ordinary LoRA, merging computes $W_{\mathrm{merged}}=W+\Delta W$ once, so
the deployed linear layer has the original shape.

\begin{table}[htbp]
\centering
\caption{Trade-offs between separate and merged LoRA deployment.}
\label{tab:lora-deployment-tradeoffs}
\begin{tblr}{
    width=\textwidth,
    colspec={Q[l,0.19\textwidth] X[l] X[l]},
    row{1}={font=\bfseries, bg=\maincolor!15!white},
    hlines,
    vlines}
Concern & Separate base plus adapter & Merged model \\
Storage for many tasks & One shared base plus small adapters & One full
checkpoint per merged task \\
Task switching & Load or activate another adapter & Load another complete
model \\
Disable or inspect adaptation & Straightforward & Original factors may no
longer be available \\
Runtime compatibility & Requires adapter support and correct base revision &
Looks like an ordinary standalone model \\
Inference overhead & Implementation dependent; measure it & No adapter branch
after a correct merge \\
Quantisation workflow & Depends on supported training/inference path & Merge
and quantisation order requires validation \\
Provenance & Base and adapter versions must remain paired & One artifact, but
its derivation must still be recorded \\
\end{tblr}
\end{table}

Keep adapters separate when one base serves many tasks, rapid switching or
rollback matters, or adapter-level analysis is required.
Merge when deployment accepts only a standalone model, task switching is
unnecessary, or measured runtime simplicity matters more than shared storage.

\begin{engineernote}[Numerical equivalence is a test, not an assumption]
Before release, compare logits or deterministic outputs from separate and
merged paths on a fixed suite.
Dtype conversion, quantisation, unsupported adapter variants, target-module
mismatch, and library behaviour can prevent exact equivalence.
Define a numerical tolerance and investigate violations.
\end{engineernote}

\section{End-to-End Practical Protocol}

\begin{practicallab}[A defensible tiny-LLM LoRA study]
\textbf{Question:}
Can a rank-eight adapter improve a narrow, measurable engineering-assistant
task without unacceptable general regression or systems cost?

\textbf{Procedure:}
\begin{enumerate}
    \item archive the model card, license, configuration, tokenizer, and hashes;
    \item write intended use, prohibited use, metric, resource budget, and
          stopping rule;
    \item create and document licensed messages with source and prompt-family
          IDs;
    \item deduplicate, then split by prompt family into train, validation, and
          private test;
    \item run B0 and B1 before any training and save their raw predictions;
    \item implement rank-eight query/value LoRA and verify the
          460,800-parameter estimate;
    \item run formatting, masking, gradient, tiny-batch overfit, save, and
          reload tests;
    \item train with recorded seeds, token budget, curves, and checkpoint rule;
    \item evaluate task, robustness, safety, regression, and uncertainty
          against baselines;
    \item benchmark separate and merged inference with the fixed measurement
          matrix;
    \item write a model card, failure gallery, resource table, and
          reproducibility command.
\end{enumerate}

\textbf{Exit criterion:}
A peer can reproduce every table from archived inputs and explain whether the
adapter should be rejected, kept separate, or merged.
\end{practicallab}

\section{What This Experiment Does Not Prove}

A successful narrow adaptation does not show that the model gained general
reasoning, acquired reliable factual knowledge, became safe, or will work on
another language or domain.
An unsuccessful run does not prove that LoRA is ineffective; the data,
template, target modules, rank, optimiser, token budget, or evaluation may be
wrong.
Good science narrows these alternative explanations through baselines,
ablations, and transparent limitations.

\begin{checkpoint}
\begin{enumerate}
    \item For $d_{\mathrm{in}}=512$, $d_{\mathrm{out}}=1024$, and $r=4$,
          calculate the LoRA parameter count.
    \item Explain why a random split can leak template-generated prompts.
    \item Distinguish time to first token from decode tokens per second.
    \item Give one reason to keep an adapter separate and one reason to merge
          it.
    \item Name three conditions required to compare throughput honestly.
\end{enumerate}
\end{checkpoint}

%
\section{Deep Dive: Serving the Adapted Model}
\begin{v3topic}{Quantisation Experiment}
Compare the reference model with 8-bit and 4-bit weight variants on the same
quality, latency, throughput, and peak-memory suite.
Separate quantised loading from quantisation-aware training and record which
layers or tensors remain at higher precision.

\end{v3topic}
\begin{v3topic}{Measure the Workload Shape}
Prompt length, generated length, batch concurrency, cache reuse, and hardware
determine serving behaviour.
Report time to first token and inter-token latency in addition to total
throughput; an average can hide an unusable tail.

\end{v3topic}
\begin{v3topic}{Adapter Governance}
Version base model, tokenizer, chat template, adapter, merge procedure,
training data, licence, and evaluation result as one deployable bill of
materials.
Verify a merged artifact against the separate adapter path before promotion and
retain a rollback target.
\end{v3topic}

\chapterreview{Trace}{Adapt}{Benchmark}{Deploy}

\section{Further Reading}

The original LoRA paper develops the low-rank update and reports its empirical
trade-offs.\textsuperscript{\cite{huLoRALowRank2021}}
The SmolLM2 paper and checkpoint model card document the model family and its
training
context.\textsuperscript{\cite{benallalSmolLM2WhenSmol2025,
huggingfaceSmolLM2135MInstruct2026}}
The PEFT documentation is the operational reference for current adapter APIs
and merging behaviour.\textsuperscript{\cite{huggingfacePEFTLoRA2026}}
Biderman et al. discuss why seemingly minor evaluation conventions can change
language-model
results.\textsuperscript{\cite{bidermanLessonsTrenchesReproducible2026}}
